% Cover letter for the MDPI Energies submission.
%
% Framing: contribution-first. The paper's earlier framing led with the HVAC
% application and described the contribution as "less a new controller than a
% corrected account", which read as low novelty on editorial screening. This
% letter leads with the transferable design rule, keeps the energy relevance to
% one paragraph, and states the reproducibility case briefly.
%
% Prior desk rejections are deliberately not mentioned: they are not required
% disclosure and would invite prejudgment. The WCCM-ECCOMAS talk IS disclosed,
% because the manuscript shares an abstract with it and similarity checking will
% surface that anyway.
\documentclass[11pt,a4paper]{article}
\usepackage[T1]{fontenc}
\usepackage[utf8]{inputenc}
\usepackage{mathptmx}
\usepackage[margin=2.4cm]{geometry}
\usepackage{microtype}
\usepackage{enumitem}
\pagestyle{empty}
\setlength{\parindent}{0pt}
\setlength{\parskip}{7pt}

\begin{document}

To the Editors\\
\emph{Energies}, MDPI

Dear Editors,

We submit the article \textbf{``Surrogate Selection for Reinforcement-Learning HVAC Control:
Step Size, Not Predictive Accuracy, Predicts Training Utility''} for consideration in
\emph{Energies}, for the Special Issue \emph{``Design, Modeling, and Testing of Heating,
Ventilation, and Air Conditioning (HVAC) and Building Energy Systems for Energy Efficiency''}.

\textbf{The problem the paper solves.} Reinforcement-learning controllers for building HVAC cannot
be trained on a real building or on a high-fidelity simulator, so the field trains them against fast
learned surrogates. That forces a design decision on every such study: which surrogate should serve
as the training environment? In practice the choice is made on predictive error, and to our
knowledge no published work has tested whether that criterion actually predicts downstream control
performance. We show that it does not, and over part of the range it is inverted: on the open
BOPTEST benchmark the least accurate surrogate in our study is the only one that trains a usable
controller, the most accurate one trains a controller that fails outright, and making a surrogate
more accurate can move it from the usable regime into the failing one. We then identify a criterion that
does work, the magnitude of the surrogate's per-step state increment, and isolate it causally with a
rescaling control that holds response-surface smoothness fixed while varying only the increment.
Because the increment is a property of the surrogate alone, it can be evaluated before any
controller is trained, which makes it usable as a selection rule rather than as a post-hoc
explanation.

\textbf{Contributions.} The manuscript makes three methodological contributions:

\begin{enumerate}[leftmargin=*,itemsep=1pt,topsep=2pt]
\item a controller-independent criterion for selecting a surrogate as an RL training environment,
established against the criterion the field currently uses;
\item \emph{role-separated surrogate training}, an architecture that stops trading optimization
smoothness against physical fidelity by assigning the two requirements to two models, with the
physics-informed twin restricted to a frozen, forward-only reward censor;
\item a pre-registered falsification protocol for surrogate-based RL claims, under which four
hypotheses were fixed before the corresponding runs and two, including one of our own, are reported
as not supported.
\end{enumerate}

\textbf{Fit to the Special Issue.} The Special Issue calls for the design, modeling and testing of
HVAC and building energy systems for energy efficiency, and the manuscript addresses the modeling
and the testing halves together. It is a modeling study, in that it compares black-box, grey-box
resistance--capacitance and hybrid representations of the same zone and quantifies what each is
good for; and it is a testing study, in that it proposes and validates a test that qualifies a
model \emph{before} it is put to use, on the open BOPTEST benchmark. The central result joins the
two: the modeling choice cannot be made on modeling criteria alone, because the property that
decides control performance is invisible to the accuracy metrics by which models are normally
tested. The energy-efficiency payoff is measured rather than asserted, at comfort violation below
5\% on both evaluation windows and 85 times live-simulator throughput.

\textbf{Relevance to \emph{Energies}.} The result is directly actionable for building-energy
practice. The role-separated architecture reaches a maintenance score of 0.041 with comfort
violation below 5\% on both evaluation windows at 85 times live-simulator throughput, turning a
policy-optimization run that would take roughly 66 hours against the emulator into about 47 minutes.
More broadly, any group building a digital twin to train an HVAC controller currently qualifies that
twin on predictive error; the paper shows what to measure instead, on standard open tooling
(BOPTEST, Modelica, Stable-Baselines3) so that the check is reproducible elsewhere.

\textbf{Evidence and reproducibility.} The central comparison is replicated over three random seeds
and the usable-versus-collapse verdict is separated by more than an order of magnitude. A
matched-resolution ablation rules out model class as the explanation, and two negative controls
(warm-starting from the twin, and using it as an accuracy-oriented partner) locate the benefit in
the role assignment itself. Boundary tests across three controller families and three hydronic test
cases delimit where the recipe holds; both negative outcomes are reported. Supplementary Materials
contain the full protocol, hyperparameters, calibration stages, seed tables and per-regime
diagnostics.

\textbf{Prior dissemination.} Preliminary results were presented as a conference talk with a
published abstract at the WCCM--ECCOMAS Congress 2026 (Munich, 19--24 July 2026), under the title
``Deep reinforcement learning for energy-efficient control in smart buildings''. The present
manuscript is a substantially extended full-length study and has not been published elsewhere.

The work is original, is not under consideration by another journal, and all authors have approved
the submission. The authors declare no conflicts of interest.

Thank you for considering our manuscript.

Sincerely,

Almaz Sapargali, on behalf of all authors\\
Corresponding author\\
almaz.sapargali2001@gmail.com

\end{document}
